\clearpage
{\large\bfseries Supplementary Notes}

The six notes give the implementation details that the main text summarises in Sections 4.1, 5.1, 5.2 (i), 5.3 (iv) and 5.4; the numbers quoted are those of the archived replication-level result files.

\textbf{Supplementary Note 1. Sampler settings, initial values and convergence criteria (Section 4.1).}
Posteriors are sampled with the No-U-Turn sampler in NumPyro, four chains of 400 warm-up and 500 sampling iterations (2,000 draws), target acceptance probability 0.95 and maximum tree depth 10; the random-walk and hierarchical terms use non-centred parameterisations. The trend is sampled in exactly reparameterised coordinates---the least-squares level and slope of the group's trend path over the observed months in place of the initial level and slope $(m_1^g, \delta_1^g)$, a linear map with unit Jacobian under which the prior (17) of the main text is evaluated at the implied values---which removes the ridge between the initial slope and the slope innovations along which the sampler otherwise needs several hundred leapfrog steps per iteration. Chains start at the prior medians, except $\sigma_\eta$ and $\tau_\zeta$, which start at 0.03 and 0.01 (the scale of their posteriors), all with small jitter; the trend starts at the least-squares line through the group's log monthly rates. Convergence requires rank-normalised $\widehat{R} < 1.01$ on every parameter except the auxiliary variables---the standard-normal variables of the non-centred parameterisations, whose functionals (the trend paths $m_t^g$ and $\delta_t^g$, $\log C_i$ and $\alpha'_g$) are checked instead, the horseshoe local scales and the latent fill-in values of missing indicators---with the trend paths of a rolling-origin fit checked over the observed months only; bulk and tail effective sample sizes above 400 on the reported quantities ($\boldsymbol{\beta}$, $\alpha'$, $\mathbf{w}$, $\kappa_w$, $\log C$ and the trend paths); and no divergent transitions. A fit that fails any criterion is repeated once on the same data with 1,000 warm-up and 1,000 sampling iterations, at target acceptance 0.99 if the first fit had divergent transitions and 0.95 otherwise; a fit that still fails is recorded as failed and never replaced by a new draw of the data, and the share of failed fits is reported with the results (Table S12). Posterior predictive checks compare four statistics with their replicate distributions: the proportion of zero counts, the maximum count, the index of dispersion by group, and the lag-1 to lag-3 autocorrelation of counts. Structural comparisons (dispersion, trend) use PSIS-LOO cross-validation; predictive comparisons use leave-future-out cross-validation or the rolling-origin protocol of Section 5.3, because PSIS-LOO does not measure forecast performance in a time series.

\textbf{Supplementary Note 2. Baseline specifications (Section 5.1).}
\emph{Sigma-rule.} The plug-in SRI is the SRI at the posterior means of $\mu$ and $\alpha'$ from the M2 filter, with the 24 months before the first decision taken from the smoothed fit at the first origin; limits are set at $-1\sigma$, $-1.5\sigma$ and $-2\sigma$ below the 24-month window mean.
\emph{Control charts.} The EWMA uses $\lambda = 0.2$ and the Poisson CUSUM the reference value of Lucas for a 1.5-fold shift, both with in-control means frozen at months 13--36 and limits set for a nominal Poisson ARL$_0$ of 24 months. The CUSUM designed under the negative binomial uses the negative-binomial likelihood-ratio reference value for the same 1.5-fold shift,
\[
k \;=\; \hat{\alpha}\,\log\!\left[\frac{\hat{\alpha} + 1.5\hat{\mu}}{\hat{\alpha} + \hat{\mu}}\right] \Big/ \log\!\left[\frac{1.5\,(\hat{\alpha} + \hat{\mu})}{\hat{\alpha} + 1.5\hat{\mu}}\right],
\]
with the dispersion $\hat{\alpha}$ estimated in phase I by the pooled moment estimator within the item's group, and its limit is the $(1 - 1/24)$ quantile of the statistic's in-control distribution simulated from the fitted negative binomial---the same nominal run length, imposed under the right distribution. An EWMA designed in the same way is reported in Table S14.
\emph{Machine-learning baselines.} XGBoost uses depth 3 or 6, learning rate 0.1 and up to 300 rounds; the LSTM and the temporal fusion transformer have 16 hidden units on 12-month windows and are fitted by Adam on minibatches of 256 sequences. All three are trained on the same $\mathbf{x}_{it}$, on $\mathbf{z}$ and on the item's own count history (XGBoost: the six most recent monthly counts and the twelve-month total; the networks: the twelve-month sequence of log counts), with the same rolling-origin splits; the tree depth, the number of boosting rounds and the early-stopping point of the networks are chosen on the last twelve labelled months of each origin's training range (inner validation only).
\emph{Farrington-type threshold.} At every decision month and for every item a quasi-Poisson regression with an intercept and a linear trend (retained when significant at the 5\% level) is fitted to the preceding 36 months, past outliers are down-weighted once through their Anscombe residuals, and the estimated dispersion turns the fitted mean into a negative-binomial predictive distribution; the item alarms when the current month's count exceeds the $(1 - 1/24)$ quantile of that predictive---a nominal in-control run length of 24 months---without seasonal terms (the simulated counts have none) and without a persistence requirement (cell-by-cell results in Table S9).

\textbf{Supplementary Note 3. The pooling comparison: pseudo-true dispersion under complete pooling (Section 5.2 (i)).}
An omitted item effect $u_i \sim N(0, \tau_C^2)$ multiplies the gamma variable in (2) of the main text, and the pseudo-true dispersion---the value that maximises the expected log-likelihood of the misspecified model under the true process, with its mean free---is 1.61 at $\tau_C = 0.3$ and 1.00 at $\tau_C = 0.6$ for the item means of the reference cell (matching the variance of the product instead gives 1.56 at $\tau_C = 0.3$). The complete-pooling fits approach this limit as the information grows, from an upward bias at low information---that of the posterior median of a negative-binomial dispersion when counts are few---to 15--17\% below the truth at $I = 30$ and $\bar{\mu} \ge 0.2$ (Table S5); at $\bar{\mu} = 0.5$ and $I = 30$ their $\alpha'$ intervals have a log-scale width of 0.25 against 0.39 for the hierarchical model on the same data and cover the truth in 20\% of cases. Re-drawing five replications of the reference cell with $\tau_C = 0$ separates the two effects: correctly specified complete pooling gives 2.34 for a truth of 2, the same replications at $\tau_C = 0.3$ give 1.87---a ratio of 0.80, the pseudo-true ratio $1.61/2$---and at $\tau_C = 0.6$ they give 1.27 with 90\% intervals that never cover the truth, while the hierarchical model on the same data still covers $\alpha'$ in 93\% and $\log C_i$ in 91\% of cases at $\tau_C = 0.6$ (Table S7).

\textbf{Supplementary Note 4. Refit schedule (Section 5.3 (iv)).}
Refitting M2 every 24 months instead of six (the schedule of the model-based baselines) on the same 65 replications gives AUC-PR 0.600 against 0.602 (paired difference $-0.002$, standard deviation 0.004), a stable-window ARL$_0$ of 32 months (90\% interval 21--60) against 35, an ARL$_1$ of 20.6 against 20.8 months, and a Brier score, calibration error, detection shares, false-alarm rate, conformal coverage and set size that differ by at most 0.005, at a quarter of the fitting time (185 against 773~s per replication) and with the same share of failed origin fits (Tables S10.1--S10.2).

\textbf{Supplementary Note 5. Sensitivity to priors and to indicator information (Section 5.4).}
The priors of Section 4.1 were varied in the reference cell, five replications each, on the data of the hierarchical fits (Table S6). Centring the dispersion hyperprior ($m_\alpha$ in (16)) at $\log 1$ or $\log 5$ instead of $\log 2$ moves the posterior median of $\alpha'$ from 2.15 to 2.12 and 2.19; doubling or quintupling the half-normal scales of the trend innovations in (17) leaves the posterior medians of $\sigma_\eta$ and $\tau_\zeta$ within 0.003 of their reference values and the level path unchanged; and an AHP prior pointing the wrong way ($\mathbf{w}^{\mathrm{AHP}} = (0.375, 0.625)$) pulls the posterior median of $\mathbf{w}$ from $(0.53, 0.31)$ to $(0.47, 0.39)$ for a truth of $(0.5, 0.3)$, with coverage of $\mathbf{w}$ falling from nine to seven of ten cases and everything else essentially unchanged: the exposure prior is informative at this sample size ($s_w = 0.1$ for a consistent matrix), and a misleading matrix shifts the weights by 0.06--0.08 without touching the quantities that drive the alarms. Removing the indicator information from the truth ($\boldsymbol{\beta} = 0$ in the reference design, five replications of the robustness protocol): the model recovers $\boldsymbol{\beta}$ at zero, and the alarm pipeline on these data gives AUC-PR 0.80 against 0.81 in the reference cell at 0.31 against 0.39 false alarms per item-year, while the realised ARL$_0$ of the $q$ channel rises from 51 to 180 months because the level path is smoother without indicator-driven variation and the bootstrap threshold is correspondingly conservative (Table S8). In this design the indicator signal is smooth (AR(1) coefficient 0.8, one period ahead of the intensity) and the group trend absorbs it.

\textbf{Supplementary Note 6. Computational cost (Section 5.4).}
All fits were run on CPUs, one single-threaded process per fit with the four chains in sequence (NumPyro/JAX; the archived code also vectorises replications and chains for GPUs). With the reparameterised trend a fit of the reference size ($I = 30$, $T = 120$) takes on average 55 leapfrog steps per iteration and a median of 48~s on one core, against 511 steps and about 20~min per fit in the direct parameterisation at target acceptance 0.99. Of all 738 Study A fits, 6.1\% still failed a convergence criterion after the refit. The MCMC fits took 20 core-hours in Study A, 22 in Study B (14 for M2) and 10 in the robustness study; with the particle filters, the bootstrap calibration and the baselines, the whole programme ran in about 25 hours of wall time on two cloud cores (Intel Xeon, 2.8~GHz) and four Apple-silicon cores, excluding the sensitivity runs. Fit diagnostics, failure shares by reason and times by study and model are tabulated in Table S12.
